\begin{frame}{Wyniki potwierdzają osiągnięcie celu}
  \begin{center}
    \metric{98\%}{skuteczność rozwiązania}\hfill
    \metric{35 ms}{czas odpowiedzi}\hfill
    \metric{12\%}{poprawa względem wariantu bazowego}
  \end{center}
  \vspace{7mm}
  \imageplaceholder[0.25\textheight]{Miejsce na wykres z jednostkami, legendą i podpisanymi osiami}
\end{frame}

\begin{frame}{Porównanie wariantów}
  \centering
  \small
  \begin{tabularx}{\textwidth}{Xrrr}
    \toprule
    \textbf{Wariant} & \textbf{Jakość} & \textbf{Czas} & \textbf{Koszt} \\
    \midrule
    Bazowy & 0,82 & 52 ms & niski \\
    Proponowany & \textbf{0,94} & \textbf{35 ms} & średni \\
    Referencyjny & 0,91 & 41 ms & wysoki \\
    \bottomrule
  \end{tabularx}
  \vspace{6mm}
  \begin{block}{Interpretacja}
    Nie poprzestawaj na tabeli: dodaj jedno zdanie wyjaśniające, co wynik zmienia w ocenie rozwiązania.
  \end{block}
\end{frame}

\begin{frame}{Wnioski i dalsze prace}
  \begin{columns}[T,onlytextwidth]
    \begin{column}{0.48\textwidth}
      \textbf{Najważniejsze wnioski}
      \begin{itemize}
        \item Cel pracy został osiągnięty w zdefiniowanym zakresie.
        \item Największy wpływ na wynik miał wskazany element.
        \item Rozwiązanie spełnia przyjęte kryteria.
      \end{itemize}
    \end{column}
    \begin{column}{0.48\textwidth}
      \textbf{Dalszy rozwój}
      \begin{itemize}
        \item rozszerzenie zbioru danych lub testów,
        \item optymalizacja wydajności,
        \item wdrożenie w środowisku docelowym.
      \end{itemize}
    \end{column}
  \end{columns}
\end{frame}
